% !TeX root = ../master.tex

\chapter{The GAIC Shared Task}
\label{ch:gaic_task}

% - Touche @ CLEF 2026: Generalizable Argument Identification in Context
% - Organized by Feger et al. -- same group that identified the shortcut problem
% - Objective: predict Argument / No-Argument given sentence + metadata

\section{Task Definition}
\label{sec:task_definition}

% - Binary classification: Argument vs. No-Argument
% - Input: sentence + rich metadata (source text, annotation guidelines, dataset provenance)
% - Evaluation: cross-dataset transfer on held-out test splits + unseen evaluation dataset
% - Primary metric: Macro F1
% - Key emphasis: "generalize beyond lexical shortcuts" and "exploit rich context information"

\section{Datasets}
\label{sec:datasets}

% - 10 benchmark datasets, ~17k labeled sentences total
% - Domains: medical abstracts (ABSTRCT), financial (FINARG), legal (USELEC),
%   scientific (ARGUMINSCI, SCIARK), persuasive essays (PE), web discourse (IAM),
%   argument quality (AFS, ACQUA), argument extraction (AEC)
% - Table: dataset name, domain, size, class balance, publication year
% - Diversity in what "argument" means across datasets -> generalization challenge

\section{Context Availability}
\label{sec:context_availability}

% - Context provided "where available" -- non-uniform across datasets
% - Table showing per-dataset availability:
%   - Full context (paper + document context + guidelines): ABSTRCT, ARGUMINSCI, PE, USELEC
%   - Partial (paper + document context): FINARG, SCIARK
%   - Minimal (paper only): ACQUA, AEC, AFS, IAM
% - Evaluation dataset may differ in what's available -> must handle gracefully
% - This non-uniformity motivates the context condition design (C0--C4)
